\section{Conclusiones}

\subsection{Hallazgos}

\begin{enumerate}
    \item \textbf{El perfilado cuantifica lo que el caso de estudio 1 solo sospechaba}. Transponer $B$ baja la tasa de fallo de caché L1 en lecturas de 22.4\% a 3.0\% (7.4$\times$ menos fallos) y reduce un 12\% las instrucciones ejecutadas. La hipótesis de que el cuello de botella era la memoria queda respaldada por medición directa.

    \item \textbf{La tasa de fallo no depende del número de hilos} (22.8\% / 22.5\% / 22.8\% para $T=1/8/16$). Eso aísla la causa: el paralelismo no empeora el acceso a la memoria, pero tampoco lo arregla.

    \item \textbf{OpenMP es la estrategia más rápida de las cinco}: 4.54$\times$ en $N=4000$, por encima del kernel transpuesto (3.55$\times$) y de pthreads manual (3.11$\times$). La explicación de por qué no supera al transpuesto está en el hallazgo anterior: paralelizar multiplica la demanda sobre una caché que ya estaba saturada.

    \item \textbf{El speedup crece con el tamaño del problema}, como corresponde a un comportamiento sano: con problemas chicos domina el \emph{overhead}, con problemas grandes el cómputo lo amortiza.

    \item \textbf{El hardware define la escala según el tamaño del problema}. En el benchmark entre el Intel i5-12450HX y el AMD Ryzen 5 5500U, la B escala mejor en problemas chicos (6 núcleos físicos homogéneos contra la arquitectura híbrida de la A) y la A en problemas grandes. La ventaja en tiempo de la A crece con $N$ (de 1.34$\times$ a 3.59$\times$), en consonancia con su mayor jerarquía de caché.

    \item \textbf{Comparar estrategias exige igualar las condiciones de compilación}. La medición inicial con \texttt{-O0} para OpenMP y \texttt{-O2} para el resto llevó a la conclusión errónea de que OpenMP era la peor estrategia. Es un error metodológico, no de implementación.
\end{enumerate}

\subsection{Limitaciones}

\begin{itemize}
    \item El perfilado con Cachegrind se ejecutó con $N=500$ (verificado estable hasta $N=1000$), porque la simulación de caché es entre 50 y 100 veces más lenta que la ejecución nativa. La tasa de fallo es una propiedad del patrón de acceso y no debería variar con $N$, pero no se verificó para los tamaños mayores.
    \item GProf no es concluyente para la variante con OpenMP: solo instrumenta el hilo principal, por lo que el reparto de tiempo dentro de la región paralela no se puede atribuir con esta herramienta.
    \item La comparación entre máquinas no pudo controlar la versión del \emph{firmware}, la configuración de la BIOS ni el estado térmico de los equipos, factores que pueden introducir diferencias de varios puntos porcentuales.
\end{itemize}

\subsection{Trabajo futuro}

\begin{itemize}
    \item \textbf{Combinar las dos optimizaciones}: aplicar la directiva de OpenMP sobre el kernel \emph{transpuesto}. Es la hipótesis más interesante que queda abierta: la transposición reduce la demanda de caché y el paralelismo reparte el trabajo restante.
    \item \textbf{Reparto por bloques en la versión con OpenMP}: usar \texttt{collapse(2)} solo sobre el loop $j$ y particionar $i$ por bloques en el nivel externo, para recuperar la localidad que tenía la versión manual.
    \item \textbf{Comparar con más de dos máquinas}, para caracterizar el comportamiento del kernel en función de las características del hardware (núcleos, jerarquía de caché, ancho de banda de memoria).
\end{itemize}
